\section{Zigzags and Shortening}

\subsection{Shortening and Short Sticklebacks}

If $F=\zigzagTuple{f}{n}$ is a zigzag and 
if for some $i$,
$1 \leq i \leq n-1$,
there is a morphism
$s$ such that $f'_i \circ s = f_{i+1}$ so that the triangle in 
\[
\zigzagShortenable{x}{y}{f}{n}{i}{s}
\]
commutes then the zigzag 
\[
\arraycolsep=7pt
\zigzagShortened{x}{y}{f}{n}{i}{s}
\]
is said to be a \textit{shortening} of $F$ and is said to be \textit{at position i of F}.

\newcommand{\rngTightener}[2]
    {
          \rng{
                \rst{#1_{#2+1}} \circ #1'_{#2}
               }            
    }

\newcommand{\rngTightenerOne}[1]
    {
          \rng{
                \rst{#1_{2}} \circ #1'_{1}
               }            
    }

\newcommand{ \rstTightener}[2]
    {
          \rst{
                #1'_{#2} \circ \rng{#1_{#2}}  
               }
    } 
\begin{notebox}[Change to the above June 24 2026]
Change the above. If $F$ is \emph{not} tight
then need the shortening to be not $f_i \circ s$ as is shown in the diagram but to be \oldt{$f_i \circ \rngTightener{f}{i} \circ s$}  
\newt{$f_i \circ \rng{f'_i} \circ s$}.
\end{notebox}

\begin{observation}
\llabel{shorteningATightZigzagObservation}
If $F$ is tight then in the shortening the
morphism $f_i \circ \rng{f'_i} \circ s$ simplifies to $f_i \circ s$.
\end{observation}
Significantly, the shortening of a tight zigzag is  tight.
It is easy to check this.


Cockett shows that if the category \catcw is a range category
satisfying the addition condition RR.5 then every tight
zigzag of morphisms in \catcw has a unique shortening. 

A stickleback which cannot be shortened is said to be
a \textit{short stickleback}.

\begin{lemma}
\llabel{secondShorteningCharacterisationLemma}
\end{lemma}
If $Z$ is a zigzag of length $2$ and $s$ is a shortening of $Z$
as shown here
\[
\zigzagTwo{x}{y}{z}
\begin{arrows}
\ncarr[15]{y1}{y2} \alabel{s}
\end{arrows}
\]
then
\begin{enumerate}[(i)]
\item 
$\shorten{\tighten{Z}}{s}{1} = \shorten{Z}{s}{1}$
\item 
$\abrst{Z} = \rst{\shorten{Z}{s}{i}}$ 
\item
$\abrng{Z} = \rng{\shorten{Z}{s}{i}}$ 
\end{enumerate}
\begin{proof}
By definition, the shortening of $Z$
is the singleton zigzag
$\tuple{z_1 \circ \rng{z'_1} \circ s}$.

On the otherhand, from  lemma \lref{howToTighteningaZigZagOfLengthTwo}
the tightening $\tighten{Z}$ is the zigzag
$\tuple{ \twoTighteningOne{z},\,\twoTighteningTwo{z},\,\twoTighteningThree{z}}$.
From observation \lref{shorteningATightZigzagObservation}
it then follows that the shortening $\shorten{\tighten{Z}}{s}{i}$
is the singleton zigzag $\tuple{\twoTighteningOne{z} \circ s}$.
\begin{enumerate}[(i)]
\item %{i}
Since $z_1 \circ \rng{\rst{z_2} \circ z'_1} \circ s$ 
simplifies to $z_1 \circ \rng{z'_1} \circ s$ as follows
\begin{align*}
 z_1 \circ \rng{\rst{z_2} \circ z'_1} \circ s
           &= z_1 \circ \rng{\rst{z'_1 \circ s} \circ z'_1} \circ s
           && \mbox{substituting $z'_1 \circ s$ for $z_2$,} \\
           &= z_1 \circ \rng{z'_1 \circ \rst{s}} \circ s 
           && \mbox{by R.4,} \\
           &= z_1 \circ \rng{z'_1} \circ \rst{s} \circ s
           && \mbox{by RR.3,} \\
           &= z_1 \circ \rng{z'_1} \circ s
           && \mbox{by R.1z} 
\end{align*}
then we have
\begin{align*}
\tighten{\shorten{Z}{s}{i}}
            &= \tuple{ z_1 \circ \rng{z'_1} \circ s } \\
            &= \tuple{
                       z_1 \circ \rng{\rst{z_2} \circ z'_1} \circ s
                     } \\
            &= \shorten{\tighten{Z}}{s}{i} &&\mbox{as required.}
\end{align*}  
\item %(ii)
We need to show that 
$\rst{z_1 \circ \rng{z'_1} \circ s}= \rst{z_1 \circ \rngTightenerOne{z}}$.
First we show that 
$\rst{\rng{z'_1} \circ s}=\rst{\rngTightenerOne{z}}$
as follows
\begin{align*}
\rst{\rng{z'_1} \circ s} 
    &= \rst{\rng{z'_1} \circ \rst{s}} && \mbox{wR.4}\\
    &= \rst{ \rng{ z'_1 \circ \rst{s} } }
                                      && \mbox{RR.3} \\
    &= \rst{ \rng{ \rst{z'_1 \circ s} \circ z'_1 } }
                                      && \mbox{R.4} \\
        &= \rst{ \rng{ \rst{z_2} \circ z'_1 } }
                        && \parbox[t]{5cm}{from the initial assumption 
                        $z'_1 \circ s = z_2$.} 
\end{align*}
Now using this we get
\begin{align*}
\rst{z_1 \circ \rng{z'_1} \circ s} 
                &=\rst{z_1 \circ \rst{\rng{z'_1} \circ s}} 
                &&\mbox{by wR.4,}\\
                &=\rst{z_1 \circ \rst{\rngTightenerOne{z}}}
                && \mbox{as shown above,}\\
                &= \rst{z_1 \circ \rngTightenerOne{z}}
                &&\mbox{by wR.4.}
\end{align*}

\item %(iii)
We need to show that 
$\rng{z_1 \circ \rng{z'_1} \circ s} 
     =\rng{\rstTightener{z}{1} \circ z_2}$ 
which we can do as follows
\begin{align*}
\rng{z_1 \circ \rng{z'_1} \circ s} &= \rng{z_1 \circ \rng{z'_1} \circ s}\\
         &= \rng{z'_1 \circ \rng{z_1} \circ s}
            && \mbox{by lemma \lref{fglemma} (vii),}\\
         &= \rng{z'_1 \circ \rst{\rng{z_1}} \circ s}
            && \mbox{by RR.1,}\\
         &= \rng{
                    \rstTightener{z}{1}
                    \circ z'_1 \circ s
                 }
            &&\mbox{by R.4,}\\
         &= \rng{\rstTightener{z}{1} \circ z_2 }
            &&\parbox[t]{5cm}{from the initial assumption  $z'_1 \circ s = z_2$.} 
\end{align*}
\end{enumerate}
\end{proof}

\subsection{Prelude to a category of short zigzags}

In what follows lemma 
\lref{deferredShorteningLemma}
describes the interaction of shortening and tightening 
and is the key to showing that there is a category of short sticklebacks. Before we get to this lemma we have a series of
subordinate lemmas that are required for its proof. 
 
\begin{lemma}
\llabel{shorteningSpecialisationLemma}
In any restriction category \catcw, 
if $f' \leq g'$ and $f \leq g$ then if  
$\nudgeup{1.5cm} % moves the triangle down to make room for arc
\downTriangleShape{x_i}{y_{i-1}}{y_i}
\downToLeft{g'}[0.6]
\downToRight{g}
\leftToRightArcing[a]{s}{30}{30} 
$ commutes 
and $\rst{f}=\rst{f'}$ then
$\downTriangleShape{x_i}{y_{i-1}}{y_i}
\downToLeft{f'}[0.6]
\downToRight{f}
\leftToRightArcing[a]{s}{30}{30} 
$ commutes.
\end{lemma}
\begin{proof}
\begin{align}
f' \circ s 
   &= \rst{f'} \circ g'\circ s 
       && \mbox{because $f' \leq g'$,}\\
   &= \rst{f} \circ g 
   && \mbox{since $\rst{f'}=\rst{f}$ 
              and $g' \circ s = g$, }\\
   &= f
   &&\mbox{since $f \leq g$}
\end{align}
\end{proof}


For convenience, 
whenever the concatenations are defined, we write
$Z \cncdot j$ for $Z \cnc \tuple{j}$ and
$i \dotcnc Z$ for $\tuple{i} \cnc Z$.

\begin{corollary}
\llabel{shorteningAtightening}
\implicitref{shorteningSpecialisationLemma}
If $F$ is a zigzag and $s$ is a shortening of $F$ at position $i$,
if $r$ is a restriction idempotent such that $F \cnc r$ is defined
then $s$ is a shortening of $\tighten{F\cnc r}$ at position $i$.
\end{corollary}

\begin{lemma}
\llabel{lemmaOfThe6thOfJuly}
If $s$ is a shortening 
of a zigzag $F$  of length $n$
then
\begin{enumerate}[(i)]
\item
if $j$ is a morphism  such that $F \cnc j$  is defined then
\begin{itemize}
\item if $s$ is a shortening at position $i$ for some $i < n$ 
then $s$ is a shortening of $F \cnc j$ 
at position $i$ and
 \[
 \shorten{F \cnc j}{s}{i} 
                       = \shorten{F}{s}{i} \cnc j,
 \] 
\item
if $s$ is a shortening at position $n$
then $s \circ j$ is a shortening of $F \cnc j$ at position $n$
and
 \[
 \shorten{F \cnc j}{s \circ j}{n} 
                       = \shorten{F}{s}{n} \cnc j,
 \] 
\end{itemize}
 \item
 if $i$ is a morphism
 such that $i \cnc F$  is defined
 then $s$ is a shortening of $i \cnc F$ and
 \[
 \shorten{i \cnc F}{s}{1}
                        = i \cnc \shorten{F}{s}{1}.
 \] 
\end{enumerate}
Note the lack of symmetry between (i) and (ii). 
\end{lemma}
\begin{proof}
Follow directly from the definition of shortening.
\end{proof}

\begin{corollary}
\llabel{corollaryOfThe6thJuly}
\implicitref{lemmaOfThe6thOfJuly}
If $Z_1$ and $Z_2$ are zigzags such that $Z_1 \cnc Z_2$ 
is defined then 
\begin{enumerate}[(i)]
\item %(i)
if $s$ is a shortening of $Z_1$
at some position $i$
then there exists a shortening $s'$ of $Z_1 \cnc Z_2$
at position $i$ such that 
\[
\shorten{Z_1 \cnc Z_2}{s'}{i} = \shorten{Z_1}{s}{i} \cnc Z_2
\]
\item %(ii)
if $s$ is a shortening of $Z_2$
at some position $i$ and $Z_1$ is of length $n$
then $s$ is a shortening  $Z_1 \cnc Z_2$
at position $n + i -1$ such that 
\[
\shorten{Z_1 \cnc Z_2}{s}{n + i -1} = Z_1 \cnc \shorten{Z_2}{s}{i}
\]
\end{enumerate}
\end{corollary}

\begin{aside}
If I talk about shortening then the context
needs to be \highlight{RR.5} range category.
\end{aside}


%zigzagIsolationTailTuple{f}{i}{n}
\newcommand{\zigzagIsolationTailTuple}[3]
{
\tuple{\rng{#1_{#2+1}},#1'_{#2+1},...#1'_{#3-1},#1_{#3}} 
}
\begin{newtt}
We use the symbol $\cnc$ to denote concatenation of zigzags. The three dots are arranged along the same rising diagonal as the morphisms in our diagrams. 

Now suppose $F$ is a zigzag of length $n$, if we denote the final element of $F$ by $f_n$ and the initial element of
$G$ as $g_1$ then in the special case that $g_1$ is a restriction 
idempotent and  $\rng{f_n}=g_1$ then the concatenation
$F \cnc G$ doesn't depend on the inital element of $G$ and the n'th
forward element of the concatenation is $f_n$. To
document this special case then we introduce a new operator.
We say that $F \cnctail G$ is defined iff 
both $F \cnc G$ is defined and  $g_1 = \rng{f_n}$.
\end{newtt}

The significance of the notation in what follows is based
on the following observation
\begin{observation}
\llabel{isolationOfshorteningObservation}

If $F$, $Z$ and $G$ are zigzags such that 
$F \cnc Z \cnctail G$ is defined, If $F$ is of length $n$
and if $s$ is a shortening of $F \cnc Z \cnctail G$ at
position $n$ then 
\[
\shorten{F \cnc Z \cnctail G}{s}{i} 
   = F \cnc \shorten{Z}{s}{i} \cnc G
\]
\end{observation}

\begin{lemma}
\llabel{isolationOfShortening}
If $F= \zigzagTuple{f}{n}$ is a zigzag and if $i$ indexes a zag
with $F$, i.e. if $1 \leq i \leq n$, then
 
\begin{enumerate}[(i)]
\item if we define zigzags $L$ of length $i$, $Z$ of length $2$ and $R$ of length $n-i-1$ as follows
\begin{align*}
L &= \zigzagTuple{f}{i} \\
Z &= \tuple {\rng{f_i}, f'_i, f_{i+1}} \\
R &= \zigzagIsolationTailTuple{f}{i}{n}
\end{align*}
then $L \cnc Z \cnctail R$ is defined and
\[
F = L \cnc Z \cnctail R
\]
\item if $s$ is a shortening of $F$ at position $i$
then $s$ is a shortening of $Z$ at position $1$ 
and
\[
\shorten{F}{s}{i} = L \cnc \shorten{Z}{s}{1} \cnc R
\]
\end{enumerate}
\end{lemma}
\begin{proof}
\begin{enumerate}[(i)]
\item
It follows directly (from the definitions) that left hand side and right hand side agree except possibly in 
their i'th forward facing elements.
The i'th forward element in the left side, $\shorten{F}{s}{i}$,
is defined to be 
$f_i \circ \rng{f'_i} \circ s$. 
On the other hand the 
corresponding element of the right hand side is
$L \cnc \tuple{\rng{f_i} \circ \rng{f'_i} \circ s} \cnc R$ which
is $f_i \circ (\rng{f_i} \circ \rng{f'_i} \circ s) \circ \rng{f_{i+1}}$
which, after removing parentheses, we can  show to be identical as follows:
\begin{align*}
f_i \circ \rng{f_i} \circ \rng{f'_i} \circ s \circ \rng{f_{i+1}}
&= f_i  \circ \rng{f'_i} \circ s \circ \rng{f_{i+1}} &&\mbox{by RR.2,}\\
&= f_i \circ \rng{f'_i} \circ s \circ \rng{f'_i \circ s} 
           &&\mbox{since $f_{i+1} = f'_i \circ s$,}           \\
&= f_i \circ \rng{f'_i} \circ s \circ \rng{\rng{f'_i} \circ s} 
           &&\mbox{by RR.4,}                                  \\
&= f_i \circ \rng{f'_i} \circ s
           && \mbox{by RR.2.}
\end{align*}
\item %(ii)
This is just an application of observation \lref{isolationOfshorteningObservation}.
\end{enumerate}
\end{proof}

The concatenation of zigzags in the above lemma has a significant feature

\begin{aside}
The decomposition above can be written
\[
F = L \cnc Z \cnctail R
\]
Point (ii) above can be extracted as a separate observation
llabel {isolationOfShortening Two} along the lines that for any zigzags F, Z and R such that 
$F \cnc Z \cnctail R$ is defined. $F$ of length $i$
\[
\shorten{F \cnc Z \cnctail R}{s}{i}
     = F \cnc \shorten{Z}{s}{1} \cnc R
\] 
\end{aside}

This is the key construction used in the proof of the following:
\begin{lemma}
\llabel{TSExtractionLemma}
If $F$ is a zigzag and $s$ is a shortening of $F$ at a position $i$ then
\[
\tighten{\shorten{F}{s}{i}} = \shorten{\tighten{F}}{s}{i}
\]
\end{lemma}
\begin{proof}
\commentary{Using abbreviated notation where we write 
               such as $Z_1 \cnc \rst{Z_2}$  for what,
            strictly, should be $Z_1 \cnc \tuple{\rstdot{Z_2}}$.}
We use the construction just given so that we have
\[
F = L \cnc Z \cnc R
\]
where 
$Z$ is of length $2$ and where $L$ and $R$ are zigzags of length 
$\geq 1$, and $s$ is a shortening of $Z$.

\newcommand{\sZ}{\shorten{Z}{s}{1}}
\newcommand{\tL}{\tighten{L}}
\newcommand{\tR}{\tighten{R}}
\newcommand {\leftPartOneOfThree}
{
    \tighten{L 
             \cnc
             \rst{\sZ \cnc r}
             }
}
\newcommand {\leftPartTwoOfThree}
{
    l \cnc \sZ \cnc r
}
\newcommand {\leftPartThreeOfThree}
{
    \tighten{\rng{l \cnc \sZ}
             \cnc
             R
            }
}


We consider each side of the target identity. For the left hand side we have
\begin{align*}
\tighten{\shorten{F}{s}{i}} 
    &= \tighten{\shorten{L \cnc Z \cnc R}{s}{i}} \\
    &= \tighten{L \cnc \sZ \cnc R} 
       &&\mbox{by lemma \lref{isolationOfShortening},}\\
    &= P_1 \cnc P_2 \cnctail P_3
       && \mbox{\parbox[t]{7cm}
            {
          by lemma \lref{howToTightenAThreefoldConcatenation},
          where $\begin{aligned}[t]
                   P_1 &=  \leftPartOneOfThree\\
                   P_2 &=  \leftPartTwoOfThree\\
                   P_3 &=  \leftPartThreeOfThree \\
                       &  \text{where $l = \rng{\tL}$    
                                        and $r = \rst{\tR}$.}
                  \end{aligned}$ 
           }
              }
\end{align*}


\newcommand{\ZZL}{\tighten{Z \cnc r}}
\newcommand{\ZZ} {\tighten{l\cnc Z \cnc r }}
\newcommand{\ZZR}{\tighten{l \cnc Z}}

\newcommand{\rightPartOneOfThree}
                 {\tighten{L \cnc \rst{\ZZL}}}
\newcommand{\rightPartTwoOfThree}
                 {\shorten{\ZZ}{s \circ r }{i}}
\newcommand{\rightPartThreeOfThree}
                 {\tighten{\rng{\ZZR} \cnc R}}

Looking at the right hand side, we have 
\begin{align*}
\shorten{\tighten{F}}{s}{i}
   &= \shorten{\tighten{L \cnc Z \cnc R}}
                {s}{i} \\
   &= \shorten{Q_1 \cnc Q_2 \cnc Q_3}{s}{i} 
       && \mbox{\parbox[t]{7cm}
                {
                  by lemma 
                  \lref{howToTightenAThreefoldConcatenation},
                  where, using the idempotents $l$ and $r$
                   as defined above, we define
                   $Q_1= \rightPartOneOfThree$\\
                  $Q_2=\ZZ$\\
                  and $Q_3=\rightPartThreeOfThree$,
                }
                } \\
    &= \shorten{Q_1 \cnc Q_2 \cnctail Q_3}{s}{i} 
        && \mbox{\parbox[t]{7cm}
            {
              because  by lemma 
                  \lref{howToTightenAThreefoldConcatenation}
              $\rng{Q_2} = \rst{Q_3}$ and because 
              the initial element of $Q_3$ is a restriction
              idempotent since the initial element of $R$ is. 
            }
                }  \\
    &= Q_1 \cnc \shorten{Q_2}{s}{i} \cnc  Q_3
        && \mbox{
            by observation \lref{isolationOfshorteningObservation}.
                } 
\end{align*}

We can complete the proof by proving the three identities
corresponding to $P_1 = Q_1$,
$P_2=\shorten{Q_2}{s}{i}$ and $P_3=Q_3$.

The first of these is
\[
\leftPartOneOfThree = \rightPartOneOfThree 
\]
this follows because by lemma \lref{lemmaOfThe6thOfJuly}(i):
\[
\shorten{Z}{s}{1} \cnc r = \shorten{Z \cnc r}{s \circ r}{1} 
\]
and because by lemma \lref{secondShorteningCharacterisationLemma}(ii)
\[
\rst{\shorten{Z \cnc r}{s \circ r}{1}} =  \rst{\tighten{Z \cnc r}} 
\]
\workt{CHECK THIS}\workt{\raisebox{0.5cm}{$\Uparrow$}}.

The second of these identities is 
\[
\leftPartTwoOfThree  = \shorten{\ZZ}{s}{i}.
\]
To see that this is the case, consider the right hand side,
$\shorten{\ZZ}{s}{i}$, which is defined, as we have seen, 
since $s$ is a shortening of $\ZZ$. But also consider that
$s$ is a shortening of $Z$ and therefore 
$s \circ r$ is a shortening of $l \cnc Z \cnc r$ and therefore
$s \circ r$ is a shortening of $\ZZ$. 
Both $s$ and $s \circ r$ are shortenings of $\ZZ$ and because
because shortenings are unique (lemma ???) we have
shortening of $\ZZ$ by $s$ and shortening of $\ZZ$ by $s \circ r$ 
are identical. 
But by lemma \lref{lemmaOfThe6thOfJuly} this shortening by 
$s \circ r$ is identical to $\leftPartTwoOfThree$ and so we arrive at
the second identity.

The third identity we need to to prove is
\[
\leftPartThreeOfThree = \rightPartThreeOfThree 
\]
which follows because by lemma \lref{lemmaOfThe6thOfJuly}(ii):
\[
l \cnc \shorten{Z}{s}{1} = \shorten{l \cnc  Z}{s}{1} 
\]
and because by lemma \lref{secondShorteningCharacterisationLemma}(iii)
\[
\rng{\shorten{l \cnc  Z}{s}{1}} =  \rng{\tighten{l \cnc  Z}} 
\]
\end{proof}
\begin{aside}
The following lemma suggests that the normal form for a zigzag is the
shortening of the tightening.
\end{aside}
\begin{lemma}
\llabel{deferredShorteningLemma}
If $Z_1$, $Z_2$ are tight zigzags in
\newt{an RR.5 range category} such that the composition $Z_1 \cnc Z_2$ is defined then
\begin{enumerate}[(i)]
\item 
$shorten(tighten(shorten(Z_1) \cnc Z_2))= shorten(tighten(Z_1 \cnc Z_2)).$
\item
$shorten(tighten(Z_1 \cnc shorten(Z_2)))= shorten(tighten(Z_1 \cnc Z_2)).$
\end{enumerate}
\end{lemma}
\begin{proof}
In the left hand sides of both (i) and (ii) 
the innermost shortening is a combination of individual shortenings
and it follows by  corollary \lref{corollaryOfThe6thJuly}
and by lemma \lref{TSExtractionLemma}
that each of these individual shortenings can be moved outside
of the tightening. 
From there they are effected by the outermost shortening. 
Hence result.
\end{proof}

\subsection{The category of short sticklebacks}

If \catcw is an RR.5 range category then
 there is a category of short sticklebacks over \catc, which we denote
 $\Stk$.
Composition of short sticklebacks is defined by
\[
F \circ G = shorten(T(F\cnc G)) 
\]
Composition, so defined, is associative because:
\begin{align*}
F \circ (G \circ H)
  &= shorten(T(F \cnc shorten(T(G \cnc H)) && \mbox{definition $\circ$,} \\
  &= shorten (T(F\cnc T(G\cnc H))) && \mbox{by lemma \lref{deferredShorteningLemma} (i),} \\
  &= shorten (T(F\cnc G\cnc H))    && \mbox{\parbox[t]{5cm}{see proof of associativity of composition of tight zigzags, section \lref{categorytightzigzags}.}}
\end{align*}  
and because, likewise, 
using lemma \lref{deferredShorteningLemma} (ii),
we have also that $(F \circ G) \circ H=shorten (T(F \cnc G \cnc H))$.

\subsection{Sequential Shortening}
\begin{lemma}
\llabel{sequentialShorteningLemma}
If $F=\zigzagTuple{f}{n}:x_1 \morph y_n$ 
is a tight stickleback in \highlight{an RR.5 range category} \catcw
and if $F$ shortens to a zigzag $\tuple{f}$ for some
morphism $f:x_1 \morph y_n$ in \catcw then  there
is a sequence of morphisms $s_1,...s_{n-1}$,
$s_i: y_i \morph y_n$ in \catcw
such that all diagrams in 
\commentary{Improve the diagramming at a later date}
\[
%\arraycolsep=12pt
\zigzagSequentialShortening{x}{y}{f}{n}{x}{g}{h}{s}
\]

commute i.e. such that
\[
f_1 \circ s_1 =f
\]
and \foreachi[n-2], 
\[
f'_i \circ s_i =  f_{i+1} \circ s_{i+1}
\] 
and
\[
f'_{n-1} \circ s_{n-1} = f_n 
\]
\end{lemma}
\begin{proof}
For any series of shortenings of $F$ to a single morphism one of this set must shorten the triangle the final zag in F i.e. 
there must be a morphism $s_{n-1}:y_{n-1} \morph y_n$
such that the triangle 
$\nudgeup{1.5cm}\downTriangleShape{x_n}{y_{n-1}}{y_n}
\downToLeft{g'}[0.6]
\downToRight{g}
\leftToRightArcing[a]{s}{30}{30} 
$ 
commutes. 
Since shortenings can be applied in any order this shortening can be applied first and the resulting zigzag 
$\tuple{f_1,f'_1,...f_{n-2}, f_{n-1} \circ s_{n-1}}$
must then shorten to the morphism $f$. 
Repeating the argument we get a morphism $s_{n-2}$
such that 
$
\begin{array}{c  c c}
&\Ob{y}{n}&               \\[0.5cm]
\Obpp{y}{n} && \Obp{y}{n} \\[0.5cm]
& \Obp{x}{n}
\end{array}
\begin{arrows}
\ncarr{ynpp}{yn}\alabel{s_{n-2}}
\ncarr{ynp}{yn}\blabel{s_{n-1}}
\ncarr{xnp}{ynpp}\alabel{f'_{n-2}}
\ncarr{xnp}{ynp}\blabel{f_{n-1}}
\end{arrows}
$
commutes and shorten $F$ to a zigzag
$\tuple{f_1,f'_1,...f_{n-3}, f_{n-2} \circ s_{n-2}}$
By iteration we arrive
at morphisms $s_1,...s_{n-1}$ which shorten $F$ to a singleton (and therefore short) zigzag $\tuple{f_1 \circ s_1}$. But since the shortening of a tight zigzag is unique 
we must have $f_1 \circ s_1=f$.
\end{proof}